\documentclass[11pt,a4paper]{article}
\usepackage[utf8]{inputenc}
\usepackage{graphicx}
\usepackage{enumitem}
\usepackage{amsmath,amssymb}
\usepackage[table]{xcolor}
\usepackage{array}
\usepackage[left=15mm,right=15mm,top=20mm,bottom=25mm]{geometry}

\newcommand{\R}{\mathbb{R}}
\newcommand{\N}{\mathbb{N}}
\graphicspath{{C:/Users/teoso/Documents/Template Math Depart/}}
\setlist[enumerate,1]{label=\arabic*.,leftmargin=*,itemsep=7pt,parsep=0pt,topsep=5pt}
\setlist[enumerate,2]{label=(\alph*),leftmargin=*,itemsep=4pt,parsep=0pt,topsep=4pt}

\begin{document}
\setlength{\abovedisplayskip}{6pt}
\setlength{\belowdisplayskip}{6pt}
\setlength{\abovedisplayshortskip}{4pt}
\setlength{\belowdisplayshortskip}{4pt}

\noindent
\begin{minipage}[t]{0.12\textwidth}
  \vspace{0pt}
  \includegraphics[width=1.6cm]{ITS.png}
\end{minipage}%
\begin{minipage}[t]{0.76\textwidth}
  \vspace{0pt}
  \centering
  \textbf{EVALUASI TENGAH SEMESTER GASAL 2025/2026}\\
  \textbf{DEPARTEMEN MATEMATIKA FSAD ITS}\\
  \textbf{PROGRAM SARJANA}
\end{minipage}%
\begin{minipage}[t]{0.12\textwidth}
  \vspace{0pt}
  \centering
  \includegraphics[width=1.6cm]{M.png}
\end{minipage}

\vspace{0.4cm}
\begin{center}
  \begin{tabular}{@{}l c l@{}}
    Matakuliah    & : & Analisis 1                                \\
    Hari, Tanggal & : & Senin, 13 Oktober 2025                    \\
    Waktu / Sifat & : & 100 menit / \textit{Closed Book}          \\
    Kelas, Dosen  & : & A.\quad Dr. Sumarsini, M.Si.              \\
                  &   & B.\quad Dr. Mahmud Yunus, M.Si.           \\
                  &   & C.\quad Dr. Rinurwati, M.Si.              \\
                  &   & D.\quad Dr. Helisyah Nur Fadhilah, M.Mat.
  \end{tabular}
\end{center}

\noindent
\begin{tabular}{|p{\dimexpr\textwidth-2\tabcolsep-2\arrayrulewidth\relax}|}
  \hline
  ETS mengukur Capaian Pembelajaran yang harus dicapai dalam mata kuliah ini, yaitu: \\
  \hline
\end{tabular}

\noindent
\begin{tabular}{|>{\centering\arraybackslash}p{1.2cm}|p{\dimexpr\textwidth-4.4cm-8\tabcolsep-5\arrayrulewidth\relax}|>{\centering\arraybackslash}p{1.4cm}|>{\centering\arraybackslash}p{1.8cm}|}
  \hline
  \rowcolor{lightgray} CPMK & \multicolumn{1}{c|}{Capaian Pembelajaran Mata Kuliah}
                            & CPL                                                              & \begin{tabular}{@{}c@{}}Bobot\\CPMK\end{tabular} \\
  \hline
  1                         & Mahasiswa mampu menjelaskan Sistem Bilangan Real.
                            & 1,4,8,9                                                          & 10\%                                             \\
  \hline
  2                         & Mahasiswa mampu menjelaskan dan membuktikan konvergensi barisan,
  barisan terbatas dan barisan Cauchy.
                            & 1,4,8,9                                                          & 15\%                                             \\
  \hline
\end{tabular}

\begin{center}
  \begin{tabular}{|c|c|c|c|}
    \hline
    Soal & Bobot Soal & Bobot CPMK-1 & Bobot CPMK-2 \\
    \hline
    1    & 20\%       & 5\%          &              \\ \hline
    2    & 20\%       & 5\%          &              \\ \hline
    3    & 20\%       &              & 5\%          \\ \hline
    4    & 20\%       &              & 5\%          \\ \hline
    5    & 20\%       &              & 5\%          \\ \hline
  \end{tabular}
\end{center}

\noindent Soal:
\begin{enumerate}
  \item Buktikan dengan Induksi Matematika, bahwa:
        \[
          1+3+5+\cdots+(2n-1)=n^2,
        \]
        berlaku untuk setiap $n\in\N$.

  \item Tunjukkan bahwa jika $a,b\in\R$, maka
        \begin{enumerate}
          \item $(-a)\cdot(-b)=a\cdot b$.
          \item $-\left(\dfrac{a}{b}\right)=\dfrac{-a}{b}$ dengan $b\ne0$.
        \end{enumerate}
        (Berikan penjelasan pada tiap langkah pengerjaan.)

  \item Diketahui $S\subset\R$ suatu himpunan terbatas dan
        $aS:=\{as:s\in S\}$ dengan $a<0$.
        Buktikan bahwa:
        \begin{enumerate}
          \item $\sup(aS)=a\inf(S)$.
          \item $\inf(aS)=a\sup(S)$.
        \end{enumerate}

  \item Dengan menggunakan definisi limit barisan, tunjukkan bahwa
        \[
          \lim_{n\to\infty}\left(\frac{1}{n}-\frac{3}{n+1}\right)=0.
        \]

  \item Buktikan bahwa jika barisan $(x_n)$ konvergen ke suatu bilangan real positif $x$,
        maka ada bilangan asli $M$ sedemikian sehingga $x_n>0$ untuk semua $n\ge M$.
\end{enumerate}

\newpage
\begin{center}
  \textbf{SOLUSI}
\end{center}
\begin{enumerate}[itemsep=12pt,parsep=5pt]
  \item Untuk $n=1$, ruas kiri adalah $1$ dan ruas kanan adalah $1^2=1$.
        Jadi, pernyataan benar untuk $n=1$.

        Misalkan pernyataan benar untuk suatu $k\in\N$.
        Artinya,
        \[
          1+3+5+\cdots+(2k-1)=k^2.
        \]
        Akan ditunjukkan bahwa pernyataan juga benar untuk $n=k+1$.
        Suku berikutnya adalah $2(k+1)-1=2k+1$.
        Dengan hipotesis induksi,
        \begin{align*}
          1+3+\cdots+(2k-1)+(2k+1)
          =k^2+(2k+1)
          =k^2+2k+1
          =(k+1)^2.
        \end{align*}
        Jadi, kebenaran untuk $n=k$ mengakibatkan kebenaran untuk $n=k+1$.
        Berdasarkan induksi matematika, rumus berlaku untuk setiap $n\in\N$.

  \item Mula-mula, tunjukkan bahwa $0\cdot c=0$ untuk setiap $c\in\R$.
        Sifat distributif memberikan
        \[
          0\cdot c=(0+0)c=0\cdot c+0\cdot c.
        \]
        Tambahkan invers aditif dari $0\cdot c$ ke kedua ruas.
        Diperoleh $0=0\cdot c$.
        Sifat komutatif juga memberikan $c\cdot0=0$.

        \begin{enumerate}
          \item Karena $a+(-a)=0$, sifat distributif memberikan
                \[
                  0=0\cdot b=(a+(-a))b=ab+(-a)b.
                \]
                Tambahkan $-(ab)$ ke kedua ruas.
                Diperoleh
                \[
                  (-a)b=-(ab).
                \]

                Selanjutnya, gunakan $b+(-b)=0$.
                Dengan sifat distributif dan hasil sebelumnya,
                \begin{align*}
                  0
                   & =(-a)\cdot0
                  =(-a)(b+(-b))
                  =(-a)b+(-a)(-b)
                  =-(ab)+(-a)(-b).
                \end{align*}
                Tambahkan $ab$ ke kedua ruas.
                Diperoleh
                \[
                  (-a)(-b)=ab.
                \]

          \item Karena $b\ne0$, terdapat invers perkalian $b^{-1}$.
                Berdasarkan definisi pembagian,
                \[
                  \frac{a}{b}=ab^{-1},\qquad
                  \frac{-a}{b}=(-a)b^{-1}.
                \]
                Sifat distributif memberikan
                \begin{align*}
                  \frac{a}{b}+\frac{-a}{b} =ab^{-1}+(-a)b^{-1} =(a+(-a))b^{-1}   =0\cdot b^{-1}   =0.
                \end{align*}
                Jadi, $\dfrac{-a}{b}$ merupakan invers aditif dari $\dfrac{a}{b}$.
                Karena invers aditif tunggal,
                \[
                  -\left(\frac{a}{b}\right)=\frac{-a}{b}.
                \]
        \end{enumerate}

  \item Asumsikan $S\ne \varnothing$ dan misalkan
        \[
          m=\inf(S),\qquad M=\sup(S).
        \]
        Keduanya ada karena $S$ terbatas dan tidak kosong.
        Untuk setiap $s\in S$, berlaku $m\le s\le M$.
        Karena $a<0$, perkalian dengan $a$ membalik arah pertidaksamaan:
        \[
          aM\le as\le am.
        \]
        Jadi, $aS$ juga terbatas dan tidak kosong.

        \begin{enumerate}
          \item Dari $as\le am$, diketahui bahwa $am$ merupakan batas atas $aS$.
                Akan dibuktikan bahwa $am$ adalah batas atas terkecil.

                Ambil sembarang batas atas $u$ dari $aS$.
                Untuk setiap $s\in S$, berlaku $as\le u$.
                Bagi kedua ruas dengan $a<0$:
                \[
                  s\ge\frac{u}{a}.
                \]
                Maka, $u/a$ merupakan batas bawah $S$.
                Karena $m$ adalah batas bawah terbesar,
                \[
                  \frac{u}{a}\le m.
                \]
                Kalikan kedua ruas dengan $a<0$.
                Arah pertidaksamaan kembali berbalik, sehingga $u\ge am$.

                Jadi, setiap batas atas $aS$ tidak lebih kecil dari $am$.
                Dengan demikian,
                \[
                  \sup(aS)=am=a\inf(S).
                \]

          \item Dari $aM\le as$, diketahui bahwa $aM$ merupakan batas bawah $aS$.
                Akan dibuktikan bahwa $aM$ adalah batas bawah terbesar.

                Ambil sembarang batas bawah $\ell$ dari $aS$.
                Untuk setiap $s\in S$, berlaku $\ell\le as$.
                Bagi kedua ruas dengan $a<0$:
                \[
                  \frac{\ell}{a}\ge s.
                \]
                Maka, $\ell/a$ merupakan batas atas $S$.
                Karena $M$ adalah batas atas terkecil,
                \[
                  M\le\frac{\ell}{a}\implies aM\ge\ell
                \]
                Jadi, setiap batas bawah $aS$ tidak lebih besar dari $aM$.
                Dengan demikian,
                \[
                  \inf(aS)=aM=a\sup(S).
                \]
        \end{enumerate}

  \item Ambil sembarang $\varepsilon>0$.
        Sifat Archimedes menjamin adanya $N\in\N$ dengan
        \[
          N>\frac{4}{\varepsilon}.
        \]
        Untuk setiap $n\ge N$, pertidaksamaan segitiga memberikan
        \begin{align*}
          \left|\left(\frac{1}{n}-\frac{3}{n+1}\right)-0\right|
           & \le\frac{1}{n}+\frac{3}{n+1}            \\
           & \le\frac{1}{n}+\frac{3}{n}              \\
           & =\frac{4}{n}\le\frac{4}{N}<\varepsilon.
        \end{align*}
        Jadi, untuk setiap $\varepsilon>0$, ada $N$ yang memenuhi definisi limit.
        Dengan demikian,
        \[
          \lim_{n\to\infty}\left(\frac{1}{n}-\frac{3}{n+1}\right)=0.
        \]

  \item Diketahui $x_n\to x$ dengan $x>0$.
        Pilih $\varepsilon=x/2$.
        Pilihan ini sah karena $x/2>0$.
        Berdasarkan definisi konvergensi, ada $M\in\N$ sehingga
        \[
          n\ge M\quad\Longrightarrow\quad |x_n-x|<\frac{x}{2}.
        \]
        Ubah pertidaksamaan nilai mutlak tersebut:
        \[
          -\frac{x}{2}<x_n-x<\frac{x}{2}.
        \]
        Tambahkan $x$ pada semua ruas.
        Diperoleh
        \[
          \frac{x}{2}<x_n<\frac{3x}{2}.
        \]
        Karena $x/2>0$, diperoleh $x_n>0$ untuk semua $n\ge M$.
\end{enumerate}

\end{document}
